\Needspace{8\baselineskip}
\section{Worker III: pole-safe creative telescoping}
\label{sec:telescoping}

\subsection{Beyond polynomial antidifferences}

Forge's polynomial recurrence worker solves identities such as
$q(n+1)-q(n)=p(n)$ for polynomial $p$~\cite{forge-structural}. That mechanism is
well suited to sums of powers. It does not, by itself, cover a summand like
$\binom nk^3$, where the parameter $n$ changes both the summand and the range of
summation.

Creative telescoping instead seeks a linear recurrence in $n$ whose summand is
a difference in $k$. The method is classical; reduction-based algorithms also
provide compact certificates for more general D-finite summands~\cite{brochet}.
For Forge, the useful deliverable is not an opaque recurrence guessed by a CAS.
It is a recurrence together with finite identities and explicit support
obligations sufficient to prove it in Lean.

The prototype specializes to
\[
 F_m(n,k)=\binom nk^m,\qquad
 S_m(n)=\sum_{k=0}^{n}F_m(n,k),\qquad n\in\N,
\]
with fixed positive integer $m$. The mathematical convention is that
$\binom nk=0$ for integer $k<0$ or $k>n$. A Lean front end can implement this
zero extension explicitly, or avoid negative indices through a separate
$k=0$ branch. It must not interpret $k-1$ as truncated natural subtraction in a
formula whose derivation assumed an integer index.

\subsection{A fully explicit bounded ansatz}

Choose an order $r\ge1$, a coefficient-degree bound $a$, and a total-degree
bound $b$ for a bivariate numerator. Search for
\[
 p_0(n),\ldots,p_r(n)\in\Q[n],\quad \deg p_j\le a,
 \qquad U(n,k)\in\Q[n,k],\quad\deg U\le b,
\]
with $p_r\ne0$, satisfying
\begin{equation}
\begin{split}
 &\sum_{j=0}^{r}p_j(n)
   \left(\prod_{t=1}^{j}(n+t)\right)^m
   \left(\prod_{t=j+1}^{r}(n+t-k)\right)^m\\
 &\hspace{1cm}=(n+r-k)^mU(n,k+1)-k^mU(n,k).
\end{split}
\label{eq:cleared-telescope}
\end{equation}
Empty products are one. The certificate equation is an identity in the
polynomial ring $\Q[n,k]$.

Write each unknown polynomial in the specified monomial basis and equate
coefficients. This is a homogeneous rational linear system with
\[
 (r+1)(a+1)+\binom{b+2}{2}
\]
unknowns. A nullspace basis describes all solutions in that ansatz. If any
solution has nonzero $p_r$, at least one nullspace basis vector does. Common
polynomial factors, rational denominators, and integer content can then be
removed consistently from all $p_j$ and $U$ to obtain a compact representative.

\par\noindent\begin{minipage}{\linewidth}
\begin{lstlisting}
for a scheduled order r and degree bounds a, b:
    introduce coefficients for p[0..r](n) and U(n,k)
    form the difference of the two sides of the polynomial identity
    M := its exact coefficient matrix
    for a nullspace basis vector with p[r] != 0:
        normalize common factors and rational content
        regularize the summation certificate
        check lower boundary and every upper support boundary
        check the leading coefficient and initial values
        return the finite certificate if all checks succeed
return UNKNOWN for the exhausted scheduled ansatz
\end{lstlisting}
\end{minipage}\par

The implementation is a bounded coefficient search, not a full implementation
of Gosper's or Zeilberger's algorithms. A failed ansatz is not evidence that no
recurrence exists. The exact linear system is complete for its chosen
coefficient space; the scheduled degree bounds and the final leading-coefficient
recognizer impose additional restrictions.

\subsection{Where the apparent poles come from}

Introduce
\[
 D(n,k)=\prod_{t=1}^{r}(n+t-k),\qquad
 A(n)=\prod_{t=1}^{r-1}(n+t).
\]
A conventional rational certificate suggested by the ansatz is
\[
 R(n,k)=\frac{k^mU(n,k)}{D(n,k)^m}.
\]
Formally, away from poles, one would use
$G(n,k)=R(n,k)F_m(n,k)$ and seek
\[
 \sum_{j=0}^{r}p_j(n)F_m(n+j,k)=G(n,k+1)-G(n,k).
\]
But $D(n,k)$ vanishes at $k=n+1,\ldots,n+r$, precisely where the shifted
summands can still be nonzero. Substituting a pole and claiming that a
zero-valued $F_m$ cancels it is invalid. In Lean, where field division is total,
blind evaluation can be even more misleading: the totalized expression is not
the removable extension used in the mathematical derivation.

The remedy is to change the certificate representation before specializing it.
Define the \emph{regularized} quantity
\begin{equation}
 H(n,k)=U(n,k)\binom{n+r-1}{k-1}^{m}.
 \label{eq:regularized-H}
\end{equation}
It is defined for every integer $k$ by the zero-extension convention. The
identity we actually prove is
\begin{equation}
 A(n)^m\sum_{j=0}^{r}p_j(n)F_m(n+j,k)
   =H(n,k+1)-H(n,k).
 \label{eq:regularized-telescope}
\end{equation}
No rational pole appears in this statement.

\subsection{Deriving the interior identity without illegal cancellation}

For $1\le k\le n$, all factors of $D(n,k)$ are positive integers. Factorial
identities give
\begin{align*}
 \frac{\binom{n+j}{k}}{\binom nk}
 &=\frac{\prod_{t=1}^{j}(n+t)}{\prod_{t=1}^{j}(n+t-k)},\\
 \frac{\binom{n+r-1}{k-1}}{\binom nk}
 &=\frac{kA(n)}{D(n,k)},\\
 \frac{\binom{n+r-1}{k}}{\binom nk}
 &=\frac{A(n)(n+r-k)}{D(n,k)}.
\end{align*}
These identities are used only where their denominators are nonzero. In
particular, the last two are applied directly to $H$; they do not evaluate
$R(n,k+1)$ at its possible pole when $k=n$.

Divide the proposed identity by the nonzero $\binom nk^m$, substitute these
ratios, and multiply by $D(n,k)^m$. The result is $A(n)^m$ times
Equation~\eqref{eq:cleared-telescope}. Since $A(n)\ne0$ for $n\ge0$, the
polynomial identity establishes the interior case. In a Lean proof, the same
argument should preferably be organized as cross-multiplied factorial lemmas,
with the positivity hypotheses explicit, rather than as a sequence of loosely
controlled rational simplifications.

\subsection{Boundary equations are part of the certificate}

The ratio argument does not cover $k=0$ or $k=n+t$ for $1\le t\le r$. These
are handled symbolically, not by evaluating finitely many numerical $n$.

At $k=0$, every $F_m(n+j,0)=1$, $H(n,0)=0$, and $H(n,1)=U(n,1)$. The required
identity is therefore
\begin{equation}
 A(n)^m\sum_{j=0}^{r}p_j(n)=U(n,1).
 \label{eq:lower-boundary}
\end{equation}
At $k=n+t$, only shifts $j\ge t$ contribute. Symmetry of binomial coefficients
reduces them to fixed-lower-index polynomials:
\[
 \binom{n+j}{n+t}=\binom{n+j}{j-t}.
\]
For a fixed nonnegative integer $q$, the expression
\[
 \binom{n+c}{q}=\frac{1}{q!}\prod_{u=0}^{q-1}(n+c-u)
\]
is a rational polynomial in $n$ on the required natural domain. The two
boundary terms from $H$ similarly involve
$\binom{n+r-1}{r-t-1}$ and $\binom{n+r-1}{r-t}$, with a negative lower index
interpreted as zero. Thus each of the $r$ upper-boundary conditions becomes a
univariate rational polynomial identity.

The independent checker explicitly expands and verifies all these identities.
Their cost is small compared with the coefficient search, and they prevent the
most obvious mismatch between generic hypergeometric algebra and the intended
finite sum.

\begin{theorem}[Pole-safe summation certificate]
Assume Equation~\eqref{eq:cleared-telescope}, the lower-boundary identity, and all
$r$ upper-boundary identities hold. Then for every $n\in\N$,
\begin{equation}
 \sum_{j=0}^{r}p_j(n)S_m(n+j)=0.
 \label{eq:sum-recurrence}
\end{equation}
\end{theorem}
\begin{proof}
The interior and boundary arguments prove Equation~\eqref{eq:regularized-telescope}
for every $k$ in the common range $0\le k\le n+r$. Sum over that range. Each
shifted summand $F_m(n+j,k)$ vanishes beyond $n+j$, so its sum is exactly
$S_m(n+j)$. The right side telescopes to
\[
 H(n,n+r+1)-H(n,0)=0,
\]
since the upper binomial coefficient has lower index $n+r$ and upper index
$n+r-1$, while the lower endpoint has lower index $-1$. The left side is
$A(n)^m$ times Equation~\eqref{eq:sum-recurrence}. Every factor of $A(n)$ is
positive, and the empty product when $r=1$ is one. Cancel this nonzero scalar.
\end{proof}

No asymptotic limit or infinite summation theorem is needed. The difficulty is
uniform handling of a moving finite support, not convergence.

\subsection{A recurrence is not yet an equality with a proposed sequence}

To characterize a sequence uniquely from an order-$r$ recurrence, require
$p_r(n)\ne0$ for every $n\ge0$ and supply the first $r$ values. If sequences $S$
and $T$ share those initial values and satisfy the same recurrence, their
difference $D$ satisfies the homogeneous recurrence. Strong induction gives
$D(n),\ldots,D(n+r-1)=0$, after which $p_r(n)D(n+r)=0$ forces $D(n+r)=0$.

The prototype accepts a simple sufficient nonvanishing certificate: either
$p_r$ or $-p_r$ has a strictly positive constant coefficient and all coefficients
nonnegative. More sophisticated examples should route this obligation to
Forge's existing univariate sign machinery rather than weaken the condition.
The initial values are computed by finite exact binomial sums.

There are consequently two different useful outputs. A \emph{recurrence proof}
establishes Equation~\eqref{eq:sum-recurrence}. A \emph{sequence-characterization
proof} adds nonvanishing and initial data. The prototype's acceptance profile
requires the latter side conditions even when only the recurrence is being
recorded. A production interface should expose these as different result
constructors.

The mutation tests multiply every $p_j$ and $U$ by $n-2$. This preserves the
interior and boundary polynomial identities and produces a true recurrence,
but its leading coefficient vanishes at $n=2$. The checker reaches and rejects
the nonvanishing step. Such a recurrence is not false; it simply does not
justify the same unrestricted uniqueness argument.

\subsection{Three synthesized certificates}

\paragraph{First powers.}
For $m=1$, choose $r=1,a=0,b=0$. The three-equation system in three unknowns has
a one-dimensional nullspace. Normalization gives
\[
 p_0=-2,\qquad p_1=1,\qquad U=-1.
\]
Thus $S_1(n+1)=2S_1(n)$ and $S_1(0)=1$. Ordinary induction gives $S_1(n)=2^n$.
The recorded artifact certifies the recurrence and its initial data; this last
closed-form identification is the additional mathematical argument just given.

\paragraph{Squares.}
For $m=2$, use $r=1,a=1,b=1$. The nine-equation system in seven unknowns yields
\[
 p_0=-2(2n+1),\qquad p_1=n+1,\qquad U=2k-3n-3.
\]
The regularized certificate is particularly small:
\[
 H(n,k)=(2k-3n-3)\binom{n}{k-1}^2.
\]
It proves
\begin{equation}
 (n+1)S_2(n+1)=2(2n+1)S_2(n),\qquad S_2(0)=1.
 \label{eq:square-recurrence}
\end{equation}
For example, the upper boundary at $k=n+1$ has left side $n+1$ and right side
$0-(-(n+1))=n+1$. The lower boundary gives $-3n-1$ on both sides.

The sequence $\binom{2n}{n}$ has the same initial value and ratio
$2(2n+1)/(n+1)$, as follows from factorials. Equation~\eqref{eq:square-recurrence}
and uniqueness therefore prove the familiar identity
$S_2(n)=\binom{2n}{n}$. The factorial ratio is a source-level proof obligation,
not a string returned by the telescoper and automatically trusted.

\paragraph{Cubes.}
For $m=3$, the executed search uses $r=2,a=2,b=5$: thirty unknowns, forty-five
coefficient equations, and a one-dimensional nullspace. It finds
\begin{align*}
 p_0(n)&=-8(n+1)^2,\\
 p_1(n)&=-7n^2-21n-16,\\
 p_2(n)&=(n+2)^2,
\end{align*}
and $U(n,k)=(n+1)^2P(n,k)$, where
\begin{align*}
 P(n,k)={}&4k^3-18k^2n-30k^2+27kn^2+93kn+78k\\
          &-14n^3-74n^2-128n-72.
\end{align*}
Here $A(n)=n+1$ and
\[
 H(n,k)=(n+1)^2P(n,k)\binom{n+1}{k-1}^{3}.
\]
The resulting Franel-number recurrence is
\begin{equation}
 (n+2)^2S_3(n+2)
 =(7n^2+21n+16)S_3(n+1)+8(n+1)^2S_3(n),
 \label{eq:franel}
\end{equation}
with $S_3(0)=1$ and $S_3(1)=2$. This is a certified recurrence, not a claim that
the algorithm found a single hypergeometric closed form for $S_3$.

All three certificates are stored in \code{results/certificates.json} with the
original summand-family parameters. The checker verifies the bivariate
identity, $r+1$ boundary rows, leading-coefficient positivity, and the initial
values. Additional tests compare the recurrences against direct finite sums
for $n=0,\ldots,30$; those numerical regressions are not substituted for the
symbolic proof obligations.

\subsection{A route to a broader summation worker}

The next implementation layer should separate summand recognition from
operator search. Recognize a source summand only after proving its shift
relations and support description. Then permit a more capable search backend
to propose an operator and certificate, while retaining the same source-level
obligations for summation.

For hypergeometric antidifferences, SymPy documents \code{gosper_normal},
\code{gosper_term}, and \code{gosper_sum}~\cite{sympy-concrete}. For D-finite
operator manipulation, Sage's \code{ore_algebra} provides creative telescoping,
closure operations, transformations, guessing, and desingularization~\cite{ore}.
The reduction-based work of Brochet and Salvy is a source of algorithms for
compact definite-summation certificates~\cite{brochet}. These are proposed
backends, not dependencies exercised by the included tests.

Operator serialization must fix a normal form. For example, write
$L=\sum_jp_j(n)E_n^j$ with coefficients on the left, where
$E_nf(n)=f(n+1)$. The rule $E_n n=(n+1)E_n$ means that ordinary commutative
polynomial multiplication is not a correct implementation of operator
composition. A recurrence checker that merely evaluates the listed coefficient
polynomials at the current $n$ avoids this ambiguity; a broader operator engine
must preserve it explicitly.

Finally, support handling should become a first-class component. A general
summand may have different moving lower and upper bounds, internal poles,
exceptional small parameters, or a nonzero boundary contribution. In those
cases the output recurrence can be inhomogeneous. The worker must return the
boundary term and prove it, not discard it because the generic certificate was
computed successfully.
